\chapter{Conclusão}
\label{cap:conclusao}

Este trabalho tratou o Problema de Alocação de Leitos (PBA) em sua versão voltada ao planejamento tático do censo hospitalar, na qual um conjunto conhecido de pacientes, com datas de admissão e de alta definidas, deve ser alocado aos quartos do hospital ao longo de um horizonte. O problema foi formulado como um modelo de Programação Linear Inteira com função objetivo dada pela soma ponderada de penalidades por especialidade inadequada, transferências e quartos mistos, e resolvido por duas abordagens. A abordagem exata aplicou o \textit{solver} Gurobi ao modelo completo, e a abordagem heurística utilizou um Algoritmo Memético, que combina um Algoritmo Genético com busca local VND e trata a capacidade dos quartos como restrição flexível na busca local, com uma etapa de reparo que mantém a população viável a cada geração. As duas abordagens foram comparadas em instâncias sintéticas de quatro portes, com dez conjuntos independentes por porte e o mesmo teto de tempo para ambas.

Os resultados mostraram dois comportamentos distintos. Na instância pequena, com 100 pacientes, o modelo exato é tratável. O Gurobi obteve a melhor solução em nove dos dez conjuntos, com \textit{gap} de 1\% a 16\%, e a heurística ficou, em média, 1,6\% acima dele e no máximo 8,2\% acima do ótimo. A partir da instância média, com 400 pacientes, a situação se inverte. O Gurobi deixa de fechar o \textit{gap} dentro do teto de tempo e, nas instâncias grande e muito grande, não chega a concluir a relaxação inicial, enquanto a heurística obtém o menor $Z$ em todos os 30 conjuntos, com redução de 55\% a 96\% em relação à solução incumbente do \textit{solver}, sempre com soluções que respeitam a capacidade dos quartos. Em nenhum dos 40 conjuntos o Gurobi comprovou a otimalidade. Esse comportamento confirma, nas instâncias deste trabalho, a dificuldade dos métodos exatos em instâncias de grande porte relatada na literatura \cite{ceschia2012modeling, demeester2010hybrid}, e decorre principalmente do tamanho do modelo exato, que chega a cerca de 28,6 milhões de variáveis na instância muito grande. A heurística, por sua vez, avaliou de 43 a 55 mil soluções por segundo em todos os portes, graças à avaliação incremental, cujo custo depende apenas da internação do paciente movido.

Essas porcentagens comparam os métodos dentro de cada porte e, como uma solução incumbente não convergida serve de referência, medem sobretudo a dificuldade do método exato, e não a qualidade da heurística. Por isso, a qualidade da heurística foi avaliada também por limites inferiores. Além do limite do \textit{solver}, este trabalho propôs um limite inferior combinatório, obtido diretamente dos dados, que conta os pacientes-dia que, por falta de leitos da especialidade adequada, ocupam necessariamente quartos de outra especialidade. Esse limite mostrou que a heurística fica a no máximo 0,3 ponto percentual do mínimo de alocações em especialidade inadequada imposto pelos dados, com desvio padrão de até 0,5 ponto percentual entre os conjuntos. Em outras palavras, quase toda a alocação em especialidade inadequada observada é consequência da distribuição dos quartos, e não da alocação, de modo que reduzi-la exigiria uma decisão estratégica sobre a composição das alas do hospital. Pelo mesmo motivo, a queda do valor de $Z$ por paciente-dia com o porte, de 26,5 para 6,2, se deve principalmente aos dados, já que acompanha a queda do mínimo estrutural. O excedente da heurística sobre esse mínimo, que depende do método, é estável entre os conjuntos e praticamente não varia com o porte, de 3,0 a 3,3 por paciente-dia, mesmo na instância muito grande, em que a heurística dispõe de menos gerações.

A heurística mostrou-se estável, com coeficiente de variação de 1,1\% a 1,9\% entre execuções de um mesmo conjunto, valor muito menor que a variação entre conjuntos, que chega a 32,3\% na instância pequena. A dispersão dos resultados se deve, portanto, principalmente às diferenças entre as instâncias, e não à aleatoriedade do método. As transferências observadas, que chegam a cerca de 15 por 100 pacientes nas instâncias maiores, não indicam uma falha do algoritmo, mas o equilíbrio definido pelos pesos, já que, com $W_{transf}/W_{spec} = 0{,}5$, transferir um paciente para o quarto adequado reduz $Z$ sempre que lhe resta ao menos um dia de internação. O mesmo raciocínio explica parte dos quartos mistos, que a função objetivo prefere quando separar os gêneros exigiria alocar pacientes fora de sua especialidade.

Conclui-se que o objetivo do trabalho foi atingido. O Algoritmo Memético fornece soluções de boa qualidade para o planejamento do censo em todos os portes avaliados, próximas do ótimo nas instâncias menores e muito superiores às obtidas pelo método exato no mesmo tempo nas instâncias maiores, que são justamente as típicas de grandes instituições hospitalares.

\section{Limitações}

A etapa de reparo, que garante a viabilidade das soluções heurísticas, tem custos. Ela consome de 8\% a 26\% do teto de tempo, conforme o porte, principalmente no reparo dos indivíduos gerados ao acaso, e, ao escolher a realocação mais barata, aumenta os quartos-dia mistos nas instâncias maiores, de 9\% a 15\% em relação à mesma heurística sem reparo. Além disso, essa comparação com a versão sem reparo foi feita a partir de uma execução anterior do protocolo, e não de execuções pareadas.

Os limites inferiores disponíveis também são fracos nas instâncias maiores. O limite combinatório considera apenas a parcela de especialidade e, como as transferências e os quartos mistos respondem por 40\% a 50\% do $Z$ da heurística nesses portes, permite afirmar apenas que o seu valor é, em média, no máximo cerca de 1,8 a 2,2 vezes o ótimo. A qualidade da heurística é, portanto, certificada de forma satisfatória apenas nas instâncias pequena e média.

O desempenho da heurística é limitado por fatores diferentes conforme o porte. Na instância pequena, a busca estagna por volta da 12ª de cerca de 21 gerações, e o tempo restante não se converte em melhora. Nas demais, o teto de tempo termina com a busca ainda em progresso, com cerca de 10 gerações na instância muito grande.

Por fim, os experimentos foram realizados em instâncias sintéticas e determinísticas, com distribuição de especialidades sorteada para os quartos, o que, na instância pequena, gera conjuntos sem nenhum quarto de alguma especialidade. Os parâmetros do Algoritmo Memético, em particular o peso $W_{cap}$, foram calibrados em um único conjunto da instância média e aplicados a todos os portes. Critérios presentes na literatura, como faixas etárias, isolamento, carga de trabalho da enfermagem e chegadas de emergência, não foram considerados.

\section{Trabalhos futuros}

As limitações observadas indicam as seguintes direções para trabalhos futuros:

\begin{itemize}
    \item \textbf{Reparo mais barato e ciente do gênero:} considerar o gênero dos ocupantes na escolha do quarto de destino do reparo, para conter o aumento de quartos mistos, e reduzir o tempo gasto no reparo dos indivíduos gerados ao acaso, por exemplo, construindo-os já cientes da capacidade.
    \item \textbf{Mecanismos contra a estagnação:} nas instâncias pequenas, avaliar estratégias de diversificação acionadas quando a busca deixa de melhorar, de modo a aproveitar as gerações restantes \cite{talbi2009metaheuristics}.
    \item \textbf{Mais gerações nas instâncias grandes:} reduzir o custo de cada geração, por exemplo, limitando o esforço da busca local, ou ampliar o teto de tempo, de modo que a busca não seja interrompida ainda em progresso.
    \item \textbf{Limites inferiores mais fortes:} desenvolver limites que considerem também as transferências e os quartos mistos, capazes de certificar a qualidade da heurística nas instâncias maiores.
    \item \textbf{Abordagens híbridas com o método exato:} utilizar a solução da heurística como ponto de partida para o Gurobi, ou o Gurobi para resolver subproblemas da heurística, como em abordagens que combinam meta-heurísticas e programação matemática \cite{belien2007building}.
    \item \textbf{Calibração dos pesos:} avaliar, com gestores hospitalares, a relação entre $W_{transf}$, $W_{gen}$ e $W_{spec}$, já que ela determina o número de transferências e de quartos mistos das soluções, e calibrar os parâmetros do algoritmo em mais de um porte e conjunto.
    \item \textbf{Cenários mais realistas:} aplicar o método a dados reais ou a geradores com distribuição de especialidades proporcional à demanda, incluir critérios adicionais, como faixas etárias e isolamento, e estender o modelo para a versão dinâmica do problema, com chegadas de emergência e incerteza no tempo de permanência \cite{ceschia2012modeling, schafer2023combining}.
\end{itemize}
